\chapter{Données, préparation et protocole expérimental}
\label{ch:data}

\section{Matrice de design}

\begin{lecture}
Une matrice de design est le tableau numérique reçu par le modèle. Chaque ligne
décrit un élève, chaque colonne une caractéristique. Une colonne supplémentaire
remplie de $1$ ajoute le même terme constant au score de tous les élèves.
\end{lecture}

Après préparation, on ajoute une coordonnée constante à chaque observation :
\[
  \xt_i=(1,x_{i1},\ldots,x_{in})^\top\in\R^{n+1}.
\]
La matrice de design empile les observations par lignes,
\[
  \X=\begin{bmatrix}\xt_1^\top\\ \vdots\\\xt_m^\top\end{bmatrix}
  \in\R^{m\times(n+1)}.
\]
Cette convention fixe toutes les dimensions ultérieures : les paramètres sont
des vecteurs colonnes de $\R^{n+1}$ et $\X\thv\in\R^m$.

Sur trois élèves décrits par deux notes standardisées, la construction devient
par exemple
\begin{figure}[htbp]
\centering
\begin{tikzpicture}[
  box/.style={draw=ink!55,rounded corners=2pt,fill=accentlight,
    inner xsep=7pt,inner ysep=5pt,align=center}]
  \node[box] (raw) {$
    \begin{array}{c|rr}
      &x_1&x_2\\\hline
      \text{élève 1}&-1{,}2&0{,}4\\
      \text{élève 2}& 0{,}3&-0{,}8\\
      \text{élève 3}& 1{,}1&0{,}6
    \end{array}$};
  \node[box,right=30mm of raw] (design) {$
    \X=\begin{bmatrix}
      1&-1{,}2& 0{,}4\\
      1& 0{,}3&-0{,}8\\
      1& 1{,}1& 0{,}6
    \end{bmatrix}$};
  \draw[-{Latex},thick] (raw)--node[above,align=center]{ajout de la\\colonne du biais}(design);
  \coordinate (midboxes) at ($(raw.south)!.5!(design.south)$);
  \node[below=4mm of midboxes,align=center]
    {une ligne $=$ une observation ; une colonne $=$ une caractéristique\\
     la première colonne, constante, est associée à $\theta_0$};
\end{tikzpicture}

\vspace{4mm}
\[
  \underbrace{\begin{bmatrix}
    1&-1{,}2&0{,}4\\1&0{,}3&-0{,}8\\1&1{,}1&0{,}6
  \end{bmatrix}}_{\X\;(3\times3)}
  \underbrace{\begin{bmatrix}\theta_0\\\theta_1\\\theta_2\end{bmatrix}}_{\thv\;(3\times1)}
  =
  \underbrace{\begin{bmatrix}
    \theta_0-1{,}2\theta_1+0{,}4\theta_2\\
    \theta_0+0{,}3\theta_1-0{,}8\theta_2\\
    \theta_0+1{,}1\theta_1+0{,}6\theta_2
  \end{bmatrix}}_{\X\thv=(z_1,z_2,z_3)^\top\;(3\times1)}.
\]
\caption{Construction concrète d'une matrice de design. La colonne de $1$
multiplie $\theta_0$ dans chaque score et encode ainsi le terme constant.
Ici $\X\in\R^{3\times3}$, $\thv\in\R^3$ et $\X\thv\in\R^3$.}
\label{fig:design-matrix-example}
\end{figure}
\FloatBarrier

\begin{procedure}{Former la matrice de design}
\textbf{Entrée :} $m$ observations déjà préparées, chacune décrite par les mêmes
$n$ caractéristiques dans un ordre fixé. \textbf{Sortie :} une matrice
$\X\in\R^{m\times(n+1)}$.
\begin{enumerate}
  \item écrire les $n$ valeurs de l'observation $i$ sur la ligne $i$ ;
  \item ajouter à gauche une colonne de $1$, associée au biais $\theta_0$ ;
  \item enregistrer l'ordre des colonnes, par exemple
        $(1,\text{Astronomy},\text{Herbology},\ldots)$ ;
  \item vérifier que chaque ligne contient $n+1$ nombres et que $\X\thv$ produit
        exactement $m$ scores.
\end{enumerate}
\end{procedure}

Avec deux caractéristiques, le rôle de la constante apparaît dans l'équation de
la frontière. Le score $\theta_1x_1+\theta_2x_2$ place la frontière sur une
droite passant par l'origine. Le score $\theta_0+\theta_1x_1+\theta_2x_2$ la
place sur $\theta_0+\theta_1x_1+\theta_2x_2=0$, qui se déplace parallèlement à
elle-même quand $\theta_0$ varie. Le biais est le paramètre qui règle cette
translation.

\begin{figure}[htbp]
\centering
\begin{tikzpicture}
\begin{axis}[courseplot,axis lines=middle,xlabel={$x_1$},ylabel={$x_2$},
  xmin=-3,xmax=3,ymin=-3,ymax=3,legend style={at={(.5,1.03)},anchor=south,
  legend columns=3,draw=none}]
  \addplot[ink,very thick,domain=-3:3]{-.7*x};
  \addlegendentry{$\theta_0=0$}
  \addplot[accent,very thick,domain=-3:3]{1.3-.7*x};
  \addlegendentry{$\theta_0<0$}
  \addplot[warning,very thick,domain=-3:3]{-1.3-.7*x};
  \addlegendentry{$\theta_0>0$}
  \addplot[only marks,mark=*,ink] coordinates {(0,0)};
\end{axis}
\end{tikzpicture}
\caption{Effet du biais pour des poids $\theta_1$ et $\theta_2$ fixés. Modifier
$\theta_0$ translate la frontière parallèlement à elle-même ; sans biais, elle
est contrainte de passer par l'origine.}
\label{fig:bias-geometry}
\end{figure}

\section{Valeurs manquantes et standardisation}

\begin{lecture}
Le calcul exige un nombre dans chaque case et des échelles comparables. Imputer
consiste à définir la valeur utilisée lorsqu'une mesure manque. Standardiser
consiste à exprimer ensuite chaque mesure par rapport à la moyenne et à
l'écart-type calculés sur l'apprentissage.
\end{lecture}

Les opérations de la régression logistique portent sur des nombres. Chaque case
vide demande donc une règle explicite, par exemple l'imputation médiane, et la
valeur d'imputation de chaque variable rejoint l'état conservé du modèle. La
suppression d'une ligne relève du même principe : un choix, énoncé et documenté.

Exemple : les valeurs d'apprentissage d'un cours sont $\{4,7,9,12,100\}$. La
médiane vaut $9$ et remplace une note manquante. La moyenne vaudrait $26{,}4$,
tirée vers le haut par la seule valeur $100$. La médiane offre donc ici une
robustesse aux valeurs extrêmes ; en contrepartie, l'imputation contracte la
dispersion et absorbe un éventuel mécanisme informatif derrière l'absence.

L'ordre retenu dans ce document --- imputer, puis estimer $\mu_j$ et $s_j$ sur la
colonne complétée --- applique une règle unique à l'apprentissage, à la validation
et à toute observation future. Il contracte $s_j$ d'autant que la proportion
d'imputations est grande. Estimer $\mu_j$ et $s_j$ sur les seules valeurs
observées est l'autre convention défensible. Le point déterminant est
l'enregistrement du choix et son application identique à chaque étape.

\begin{figure}[htbp]
\centering
\begin{tabular}{c@{\hspace{14mm}$\xrightarrow{\text{médiane }=9}$\hspace{14mm}}c}
\begin{tabular}{c}\toprule valeur observée\\\midrule
$4$\\$7$\\\cellcolor{warning!25}\textcolor{warning}{manquante}\\$12$\\$100$\\\bottomrule
\end{tabular}
&
\begin{tabular}{c}\toprule valeur préparée\\\midrule
$4$\\$7$\\\cellcolor{warning!25}$9$\\$12$\\$100$\\\bottomrule
\end{tabular}
\end{tabular}
\caption{Imputation médiane sur une variable. La case colorée porte une valeur
calculée, d'un statut distinct des quatre mesures qui l'entourent. La médiane
provient des seules lignes d'apprentissage.}
\label{fig:median-imputation}
\end{figure}

Pour une note brute $r_j$, on estime sur l'apprentissage
\[
  \mu_j=\frac1{m_{\mathrm{tr}}}\sum_{i\in\mathcal I_{\mathrm{tr}}}r_{ij},
  \qquad
  s_j=\sqrt{\frac1{m_{\mathrm{tr}}}\sum_{i\in\mathcal I_{\mathrm{tr}}}
    (r_{ij}-\mu_j)^2},
  \qquad x_j=\frac{r_j-\mu_j}{s_j}.
\]
Le diviseur $m_{\mathrm{tr}}$ ou $m_{\mathrm{tr}}-1$ change le résultat de moins
d'un pour mille sur $1\,600$ lignes ; le même reste employé à l'entraînement et à
la prédiction. Un $s_j$ nul signale une colonne constante sur l'apprentissage,
que le programme retire de la liste des variables.

La standardisation reparamètre le modèle et conserve la famille des frontières
linéaires. Elle améliore le conditionnement numérique de l'optimisation et rend
une éventuelle pénalisation comparable entre variables
\cite[sec.~4.6.3]{hastie2009}.

Le \emph{conditionnement} mesure la sensibilité du calcul aux directions de
l'espace des paramètres. Une variable qui varie autour de $10^3$ et une autre
autour de $10^{-2}$ répondent à un même déplacement de leurs coefficients par des
variations de score dans un rapport de $10^5$. Les courbes de niveau du coût
s'allongent d'autant, et la descente zigzague ou progresse par pas très courts.
La figure~\ref{fig:contours} en donne la représentation géométrique. La
standardisation ramène les échelles à $1$ et conserve l'ordre des observations
dans chaque colonne.

\begin{resultat}[Exemple de standardisation]
Supposons que la note Astronomy ait, sur l'apprentissage, une moyenne
$\mu=20$ et un écart-type $s=100$. Les notes brutes $-180$, $20$ et $170$
deviennent respectivement $-2$, $0$ et $1{,}5$. La transformation conserve leur
ordre et leurs positions relatives, et fixe une origine et une échelle communes.
Une valeur standardisée de $1{,}5$ se lit « un écart-type et demi au-dessus de la
moyenne d'apprentissage », dans l'unité propre à cette colonne.
\end{resultat}

\begin{figure}[htbp]
\centering
\begin{tikzpicture}[x=1cm,y=1cm]
  \node[anchor=east] at (.4,1.25){note brute $r$};
  \node[anchor=east] at (.4,-.35){valeur $x=(r-20)/100$};
  \draw[{Latex}-{Latex},thick] (.7,1.25)--(11.3,1.25);
  \draw[{Latex}-{Latex},thick] (.7,-.35)--(11.3,-.35);
  \foreach \xpos/\raw/\std in {1/-240/-2{,}6,11/240/2{,}2}{
    \draw (\xpos,1.16)--(\xpos,1.34) node[above=2pt]{\raw};
    \draw (\xpos,-.44)--(\xpos,-.26) node[below=2pt]{\std};
  }
  \foreach \xpos/\raw/\std in {2.25/-180/-2,6.4167/20/0,9.5417/170/1{,}5}{
    \fill[accent] (\xpos,1.25) circle (2.5pt) node[above=7pt]{\raw};
    \fill[warning] (\xpos,-.35) circle (2.5pt) node[below=7pt]{\std};
    \draw[-{Latex},dashed,gridgray!80] (\xpos,1.02)--(\xpos,-.12);
  }
  \node[draw=ink!45,fill=white,rounded corners=2pt,inner sep=4pt]
    at (6,2.35){$r=20+100x$ : même position, autre origine et autre unité};
\end{tikzpicture}
\caption{Correspondance complète entre l'axe des notes brutes et l'axe
standardisé. Chaque flèche verticale relie deux écritures de la même observation :
$-180\leftrightarrow-2$, $20\leftrightarrow0$ et $170\leftrightarrow1{,}5$.}
\label{fig:standardization}
\end{figure}

\begin{vigilance}[Fuite de données]
Calculer $\mu_j$, $s_j$, une médiane d'imputation ou choisir des variables à
partir de l'ensemble complet transmet de l'information du test vers
l'apprentissage, et rend l'évaluation optimiste. La partition précède donc toute
transformation apprise des données.
\end{vigilance}

\begin{procedure}{Préparer une caractéristique numérique}
\textbf{Entrée :} les valeurs brutes $r_{ij}$ de la caractéristique $j$ sur les
lignes d'apprentissage. \textbf{Sortie :} la valeur d'imputation, $\mu_j$, $s_j$
et les valeurs transformées $x_{ij}$.
\begin{enumerate}
  \item calculer la médiane des valeurs présentes et l'utiliser pour remplacer
        les valeurs manquantes ;
  \item calculer ensuite $\mu_j$ et $s_j$ sur cette colonne d'apprentissage complétée ;
  \item transformer chaque valeur par $x_{ij}=(r_{ij}-\mu_j)/s_j$ ; par exemple,
        $r=170$, $\mu=20$, $s=100$ donnent $x=1{,}5$ ;
  \item conserver médiane, $\mu_j$ et $s_j$, et les réappliquer telles quelles à
        la validation, au test et à toute observation future ;
  \item si $s_j=0$, signaler une colonne constante et la retirer du calcul.
\end{enumerate}
\end{procedure}

\section{Partition et déséquilibre des classes}

\begin{lecture}
La partition produit deux listes d'indices disjointes : les indices
d'apprentissage $\mathcal I_{\mathrm{tr}}$ et les indices de test
$\mathcal I_{\mathrm{te}}$. Une ligne appartient à une seule liste. Toutes les
transformations et tous les coefficients sont estimés depuis la première ; la
seconde reste intacte jusqu'à l'évaluation finale.
\end{lecture}

Les deux listes vérifient
\[
  \mathcal I_{\mathrm{tr}}\cap\mathcal I_{\mathrm{te}}=\varnothing,
  \qquad
  \mathcal I_{\mathrm{tr}}\cup\mathcal I_{\mathrm{te}}=\{1,\ldots,m\}.
\]
Pour construire une partition stratifiée, on regroupe d'abord les indices par
classe. Dans chaque groupe, on mélange les indices avec une graine fixée, puis on
affecte la proportion prévue à l'apprentissage et le reste au test. Les morceaux
obtenus sont enfin réunis dans les deux listes. L'opération porte sur les indices
et laisse les données numériques en place.

L'état à conserver tient donc en trois objets : la graine et les deux listes
d'indices. Moyennes, écarts-types et valeurs d'imputation viennent ensuite des
seules lignes désignées par $\mathcal I_{\mathrm{tr}}$, puis s'appliquent telles
quelles aux lignes de $\mathcal I_{\mathrm{te}}$.

\section{Préservation des effectifs par classe}

Avec $1\,000$ élèves, dont $900$ majoritaires et $100$ dans la classe étudiée,
une répartition $80\,\%/20\,\%$ affecte séparément $80\,\%$ de chaque classe à
l'apprentissage. On obtient $720+80=800$ lignes d'apprentissage et
$180+20=200$ lignes de test. Le déséquilibre $90\,\%/10\,\%$ se retrouve à
l'identique dans les deux sous-ensembles.

\begin{figure}[htbp]
\centering
\setlength{\tabcolsep}{12pt}
\renewcommand{\arraystretch}{1.35}
\begin{tabular}{lrrr}
\toprule
classe & ensemble disponible & apprentissage ($80\,\%$) & test ($20\,\%$) \\
\midrule
\textcolor{accent}{majoritaire} & $900$ & $720$ & $180$ \\
\textcolor{warning}{classe étudiée} & $100$ & $80$ & $20$ \\
\midrule
total & $1\,000$ & $800$ & $200$ \\
\bottomrule
\end{tabular}

\vspace{.7em}
\begin{tabular}{ccc}
$900=720+180$ & \qquad & $100=80+20$
\end{tabular}
\caption{Effectifs produits par une partition stratifiée $80\,\%/20\,\%$.}
\label{fig:stratification}
\end{figure}

Un \emph{hyperparamètre} est une quantité que le protocole fixe avant
l'ajustement. Les coefficients $\thv$ sont des paramètres appris par le gradient ;
le taux $\alpha$ et la pénalisation $\lambda$ sont des hyperparamètres. Si trois
valeurs de $\lambda$ donnent sur validation des exactitudes $0{,}78$, $0{,}83$ et
$0{,}80$, la deuxième est retenue, puis le test intervient une seule fois. La
validation croisée du chapitre~\ref{ch:evaluation} étend ce principe en faisant
tourner le rôle de validation sur plusieurs blocs.

\begin{resultat}[Quantités fixées après la préparation]
La préparation produit les listes d'indices d'apprentissage et de test, une
valeur d'imputation par caractéristique, une moyenne et un écart-type par
caractéristique, puis les matrices numériques transformées. Ces quantités sont
estimées une fois sur l'apprentissage et réutilisées sans modification pour la
validation, le test et toute observation future.
\end{resultat}

\begin{procedure}{Construire une partition stratifiée}
\textbf{Entrées :} les numéros de lignes $1,\ldots,m$, leur classe $c_i$, une
proportion d'apprentissage et une \emph{graine}. La graine est un entier qui
initialise le générateur pseudo-aléatoire : reprendre le même entier reproduit
le même mélange. \textbf{Sorties :} deux listes d'indices
$\mathcal I_{\mathrm{tr}}$ et $\mathcal I_{\mathrm{te}}$.
\begin{enumerate}
  \item former une liste de numéros de lignes par classe ; par exemple, une liste
        contient les $100$ indices de la classe étudiée ;
  \item mélanger chaque liste après initialisation avec la graine enregistrée ;
  \item pour une proportion $80\,\%$, placer les $80$ premiers indices de cette
        liste dans $\mathcal I_{\mathrm{tr}}$ et les $20$ autres dans
        $\mathcal I_{\mathrm{te}}$ ; répéter pour chaque classe ;
  \item concaténer les morceaux obtenus : dans l'exemple complet, les deux listes
        contiennent respectivement $800$ et $200$ indices ;
  \item vérifier qu'aucun indice n'apparaît deux fois, qu'aucun ne manque et que
        les effectifs par classe correspondent à la table~\ref{fig:stratification} ;
  \item enregistrer définitivement graine et listes avant de calculer médianes,
        moyennes ou écarts-types.
\end{enumerate}
\end{procedure}
